\addtocontents{toc}{\protect\Needspace{3\baselineskip}}
\originalchapter{8.821 课程详讲：D 膜、规范场与几何 Higgs 机制}
本章对应第 13--15 讲的 D 膜物理, 第 15 讲末的 D3 引力解在下一章展开. $p$ 是膜的\textbf{空间}维数, 世界体有 $p+1$ 维. 用 $Y^a$ 表示有长度量纲的横向位置, 用 $\varphi^a=Y^a/(2\pi\alpha')$ 表示低能规范理论中的标量; 两种归一化从一开始分开.

\originalsection{Dirichlet 条件为什么定义一个动力学物体}
设目标坐标 $X^\alpha$, $\alpha=0,\ldots,p$, 沿膜方向满足 Neumann 条件, 横向 $X^a$, $a=p+1,\ldots,D-1$, 的端点固定在 $Y^a$. 端点可沿这张 $p$ 维面移动, 不能离开它, 这就是 D$p$ 膜. 一张固定直膜保留 $ISO(1,p)\times SO(D-p-1)$, 破坏横向平移及混合切向、法向的 Lorentz 变换. 固定膜是一种背景近似, 不意味着它在完整理论中永远不动.

时间对稳定 Lorentzian 膜通常取 Neumann 条件. D$(-1)$ 膜则在欧氏目标时空中局域于全部坐标, 是瞬子式对象, 不是一张沿真实时间持续传播的负维空间物体. 完全 Neumann 的开玻色弦可看作端点在填满空间的 D25 膜上; Type II 则有其允许的膜维数和投影, 不能把 D25 的例子照搬到十维.

同一膜两端的横向 DD 解没有 $p^a\tau$ 零模, 有
\[
X^a=Y^a+\sqrt{2\alpha'}\sum_{n\ne0}
\frac{\alpha_n^a}{n}\ee^{-\ii n\tau}\sin n\sigma.
\]
这套模式仍取 $\alpha_n^{a\dagger}=\alpha_{-n}^a$. 正弦在 $n\to-n$ 时反号, 与 $1/n$ 的反号相消, 因而此处不像 NN 余弦模式那样需要前置 $\ii$. 直接共轭是一个必要的实场检查. 动量 $p^a=0$ 不表示横向弦段从此没有振子; 正弦模式仍能在膜外振荡, 只是端点固定.

对两端同类 NN 或 DD, 频率为整数, 横向正规序常数不变. 混合 DN 情形例如 $X(0)=0,X'(\pi)=0$, 波函数取 $\sin r\sigma$, 要求 $\cos r\pi=0$, 得 $r=n+1/2$. 频率已变成半整数, 零点能也会变, 因而不能对不同维膜交叉的所有开弦都直接使用同一整数谱截距.

\originalsection{无质量矢量和位置标量的计数}
对 $p\ge1$, 可沿膜内一个空间方向选光锥. 开弦最低无质量激发分成沿膜的 $\alpha_{-1}^i$, $i=2,\ldots,p$, 和横向 $\alpha_{-1}^a$. 前者有 $p-1$ 个偏振, 正是 $p+1$ 维无质量矢量 $A_\alpha$ 的物理偏振; 后者有 $D-p-1$ 份, 是膜上的标量. 总数 $(p-1)+(D-p-1)=D-2$, 与原横向振子数一致.

一个常数横向标量把端点位置 $Y^a$ 平移. 由于平直目标中的平移对称性, 整张膜在任一常数位置都应有同样能量, 所以这些标量没有单膜平直背景中的位置势, 可看作破缺横向平移的 Goldstone 模. 它们的慢变场描述膜的弯曲与运动. 对 D0 世界体量子力学, $A_0$ 是约束场, 没有传播的矢量偏振; 不能把上面 $p-1$ 的传播态计数外推成负一个偏振.

在 Type II 的 BPS 膜上, 这些玻色场还有相应费米伙伴. 在玻色弦里, 开弦快子使同类膜构型也不稳定; 其无质量截断有形式上的作用量用途, 不等于稳定真空的全部物理. 后续 AdS/CFT 构造用的是 Type IIB 的相应稳定 D3 膜.

\originalsection{DBI 的二次展开和坐标归一化}
单张膜低导数作用量用 Dirac--Born--Infeld（DBI）描述. 在平直目标、$B=0$、常数膨胀子已吸收到张力 $T_p$ 的约定下,
\[
S_{\rm DBI}=-T_p\int\dd^{p+1}x\,
\sqrt{-\det(\eta_{\alpha\beta}+\partial_\alpha Y^a\partial_\beta Y^a
+2\pi\alpha'F_{\alpha\beta})}.
\]
更一般情形用 $P[g+B]$, 即把目标张量拉回世界体, 并有场依赖的膨胀子因子. DBI 保留合适的 $F$ 和一阶嵌入导数的高次幂, 但忽略更高导数修正; 它不是对任意剧烈变化场完全精确的作用量.

用 $\sqrt{\det(I+K)}=1+\tfrac12\operatorname{Tr}K
+\tfrac18(\operatorname{Tr}K)^2-\tfrac14\operatorname{Tr}K^2+\cdots$.
在二次场阶, $\operatorname{Tr}F=0$, 而混合指标 $F^\alpha{}_{\beta}$ 有 $\operatorname{Tr}F^2=-F_{\alpha\beta}F^{\alpha\beta}$. 因而
\[
S_{
\rm DBI}=-T_p\int\dd^{p+1}x\left[
1+\frac12\partial_\alpha Y^a\partial^\alpha Y^a
+\frac{(2\pi\alpha')^2}{4}F_{\alpha\beta}F^{\alpha\beta}+\cdots\right].
\]
标量项的 $(2\pi\alpha')$ 因子只有在用 $\varphi^a$ 替代 $Y^a$ 后才出现. $[Y]=L$、$[\varphi]=L^{-1}$、$[F]=L^{-2}$, 所以行列式内各项无量纲. 这项量纲检查可以防止把位置与规范标量混写.

若关闭 $F$、位置只随时间变化, 作用量精确变成 $-M_p\int\dd t\sqrt{1-\dot Y^2}$, $M_p=T_pV_p$, 其展开是 $-M_p+M_p\dot Y^2/2+\cdots$. 这就是整个有限体积膜质心的相对论运动, 证明 $Y^a$ 与位置解释相符, 也固定了时间动能的正确号.

\originalsection{张力为何反比于闭弦耦合}
两膜之间的 cylinder 或 annulus 振幅有两种切法: 按开放方向看是一条开弦的一环真空过程, 按封闭方向看是闭弦在两个膜间交换. Annulus 的 Euler 特征是 0, 领先振幅规模 $g_s^0$. 闭弦低能交换规模为 $G_NT_p^2$, 而 $G_N\propto g_s^2$, 匹配就要求 $T_p\propto1/g_s$. 这是从同一世界面的两个通道读出的耦合尺度, 不是把两个独立模型随意调参相等.

在 Type II 的标准膜归一化中,
\[
T_p=\frac1{(2\pi)^pg_s(\alpha')^{(p+1)/2}}.
\]
量纲为 $L^{-(p+1)}$, 与 $p$ 维空间中单位体积能量相符. 膜质量 $M_p=T_pV_p$ 在 $g_s\to0$ 时趋大, 所以膜态不是微扰闭弦 Fock 空间中普通的有限阶粒子激发. 但其上的开弦小振动仍可微扰处理. 开弦耦合由图拓扑给 $g_o^2\propto g_s$, 具体系数依场归一化.

该张力的精确常数需要 annulus 振幅与闭弦交换的明确匹配, PS3.5 分辨其耦合和长度幂次. 完整振幅及其 oscillator 迹不是由上面一行量纲论证重算出来的; 本章以已知 Type II 正规化作为后续 D3 字典的输入, 同时把尺度论证完整写出.

\originalsection{端点标签变成矩阵与 \texorpdfstring{$U(N)$}{U(N)}}
对 $N$ 张重合、同类定向膜, 开弦两端可落在任意标签 $i,j=1,\ldots,N$ 上. 一套振子态因此带 $N^2$ 个端点标签, 写 $|\Psi;ij\rangle$, 低能场也写成 $N\times N$ 矩阵. 反向弦交换两端, 对应 Hermitian 场中 $(A_{ij})^*=A_{ji}$ 等关系.

两条弦拼接要求中间端点标签相同, 所以振幅中的颜色收缩是矩阵乘法. 多条定向弦沿盘边循环拼接, 产生单迹 $\operatorname{Tr}(A_1A_2\cdots A_n)$. Chan--Paton 标签基变换 $\Psi\to U\Psi U^\dagger$ 留下这些振幅不变. 对无质量矢量, 低能一致相互作用成为世界体的 $U(N)$ Yang--Mills 规范理论, 不只是把一个全局矩阵群重命名成局部群; 这个结构由无质量矢量的规范冗余和振幅匹配共同确定.

取 $D_\alpha\varphi=\partial_\alpha\varphi-\ii[A_\alpha,\varphi]$,
$F_{\alpha\beta}=\partial_\alpha A_\beta-\partial_\beta A_\alpha-\ii[A_\alpha,A_\beta]$.
选 $\operatorname{Tr}(T^aT^b)=\delta^{ab}/2$、规范动能为 $-\operatorname{Tr}F^2/(2g_{\rm YM}^2)$ 的约定, 玻色领先作用量为
\[
S=-\frac1{g_{\rm YM}^2}\int\dd^{p+1}x\,
\operatorname{Tr}\left[
\frac12F_{\alpha\beta}F^{\alpha\beta}
+D_\alpha\varphi^aD^\alpha\varphi^a
-\frac12[\varphi^a,\varphi^b][\varphi^a,\varphi^b]
\right]+\text{费米伙伴与高阶项}.
\]
它可由高维 Yang--Mills 的降维读出: 原来法向的规范场分量成为 $\varphi^a$, 混合场强成为 $D_\alpha\varphi^a$, 两个法向场强成为 $-\ii[\varphi^a,\varphi^b]$. 对 Hermitian 标量, 对易子为反 Hermitian, 静态势 $V=-\operatorname{Tr}[\varphi^a,\varphi^b]^2/(2g^2)$ 非负, 最小值为各位置矩阵可交换. 这解释了对角位置真空, 不只是列出一个势项.

按上述迹与动能选择, 与单膜 DBI 系数比较给
\[
g_{\rm YM}^2=\frac2{T_p(2\pi\alpha')^2}
=2(2\pi)^{p-2}g_s(\alpha')^{(p-3)/2}.
\]
有的文献改用 $-\operatorname{Tr}F^2/(4g^2)$, 数值耦合相差因子 2. 本卷保持 D3 的 $g_{\rm YM}^2=4\pi g_s$, 与 AdS 字典一致. $[g_{\rm YM}^2]=L^{p-3}$, 只有 $p=3$ 这一值无长度量纲. 对多膜的全部高次 DBI, 不能简单把单膜平方根外加一个迹就称为完整答案; 非交换与导数修正已超过此处可靠的领先 Yang--Mills 截断.

\originalsection{拉伸弦质量怎样等于 Higgs 质量}
两平行膜位置为 $Y_i^a,Y_j^a$, 距离 $d_{ij}=|Y_i-Y_j|$. DD 经典解增加线性空间项 $X^a=Y_i^a+(Y_j^a-Y_i^a)\sigma/\pi+\text{正弦模}$. 它不改变振子频率, 却给世界面 Hamiltonian 一个常数伸长能量, 因而
\[
M_{ij}^2=\left(\frac{d_{ij}}{2\pi\alpha'}\right)^2
+\frac{N-a}{\alpha'}.
\]
对原本无质量的层级, 质量 $M_{ij}=T_0d_{ij}$. 膜重合 $d=0$ 时, 所有端点组合无质量, 完整 $U(N)$ 可见; 分开后非对角 $i\ne j$ 组合变重.

低能场论给同样结果. 取 $\langle\varphi^a\rangle=\operatorname{diag}(Y_1^a,\ldots,Y_N^a)/(2\pi\alpha')$. 则 $[A_\alpha,\varphi^a]_{ij}=(\varphi_j^a-\varphi_i^a)A_{\alpha,ij}$, 标量动能中的 $D\varphi$ 项成为矢量质量项, 读出
\[
M_{ij}^2=\sum_a(\varphi_i^a-\varphi_j^a)^2
=\frac{d_{ij}^2}{(2\pi\alpha')^2}.
\]
若膜分成 $k$ 个重合团、数量 $n_1,\ldots,n_k$, 未破缺群为 $\prod_{s=1}^kU(n_s)$, $\sum_sn_s=N$. 这个几何 Higgs 机制把“一个位置标量的期望”与“端点之间的一条有张力弦”定量联系起来, PS3.4 给出详细考核.

\originalsection{PS3.3 的阅读说明：边界条件与运动}
原作业 PS3.3 是一段阅读说明, 不要求另外证明一个考核小问. 它强调: 先在经典弦中把 Dirichlet 子流形当作刚性背景, 量子谱却给每个法向一个无质量位置标量, 常值标量把原边界移到另一位置, 慢变标量使它成为会弯曲和运动的 D 膜. 闭弦里类似地, 原先固定的 metric 在量子谱中出现无质量自旋二态后也可成为动态背景.

但任意背景不自动是一个一致弦背景, 还要维持世界面 Weyl 性质. 阅读说明以 D0 举例, 边界写为 $X^i(\tau,0)=X^i(\tau,\pi)=f^i(X^0)$, 即点状膜沿 $x^i=f^i(t)$ 运动. 它引用世界面一致性导出的运动条件
\[
\frac\dd{\dd t}\left[\frac{\dot f^i}{\sqrt{1-\sum_j(\dot f^j)^2}}\right]=0.
\]
从本章 DBI 的质心作用量 $-M_0\int\dd t\sqrt{1-\dot f^2}$ 作 Euler--Lagrange 变分, 正好得到同一式. 这里机械验证的是相对论粒子方程, 并非冒充已计算全部世界面 beta 函数. 原作业给 R. Leigh, \textit{Mod. Phys. Lett. A4}, 2767 (1989) 作为更一般的阅读引用, 并明确没有要求推导该 Weyl 条件; 本卷只保留这一引用, 不复录外部论文全文.

\originalsection{高次形式场为何以膜作为荷源}
以微分形式约定
$C_q=(1/q!)C_{\mu_1\cdots\mu_q}\dd X^{\mu_1}\wedge\cdots\wedge\dd X^{\mu_q}$,
$F_{q+1}=\dd C_q$. 它的规范变换是 $C_q\to C_q+\dd\Lambda_{q-1}$, 动能的分量归一化通常为 $-F_{q+1}^2/[2(q+1)!]$. 像点电荷耦合 $\int A_1$ 一样, 电性 D$p$ 膜可耦合
\[
S_{\rm WZ}=\mu_p\int_{\mathcal W_{p+1}}P[C_{p+1}].
\]
形式度数等于世界体维数, 因而可对它积分. 膜闭合或适当边界条件下, 规范变分成为边界全导数, 说明其电荷守恒. 复杂背景还有 $B,F$ 和曲率的 Wess--Zumino 耦合, 此处只取所需领先电荷项.

在 $D$ 维, 电势 $C_q$ 的磁对偶势度数为 $D-q-2$, 对应磁膜空间维数 $D-q-3$. 因而 Type IIA 的 $C_1$ 以 D0 电耦合、D6 磁耦合, $C_3$ 以 D2 电耦合、D4 磁耦合. Type IIB 的 $C_0$ 对 D$(-1)$ 与 D7, $C_2$ 对 D1 与 D5, 自对偶五形式场强的 $C_4$ 对电磁同类的 D3. 这些对应只需核对形式的次数与 Hodge 对偶, 不应把两种 Type 的名字写反.

在适当平直超弦背景, BPS D 膜保留一半超对称, 并饱和由荷决定的质量界. 守恒荷阻止完全衰变为中性态, BPS 最低能量条件提供进一步稳定性. “携最小荷”单独不排除衰变为\textbf{同荷}更低能态, 所以不能用一个荷标签代替整个稳定性论证. Bosonic 膜、膜反膜对或不满足这些背景条件的构型不自动享有同样稳定性.

对十维 IIB, 标准 Newton 常数是
\[
16\pi G_{10}=(2\pi)^7g_s^2(\alpha')^4,
\qquad G_{10}=8\pi^6g_s^2(\alpha')^4.
\]
取 D3 张力 $T_3=1/[(2\pi)^3g_s(\alpha')^2]$, 可核对 $T_3\sqrt{16\pi G_{10}}=\sqrt{2\pi}$. 这正是课程对动能标准归一化自对偶 RR 场所给的最小电磁荷. 在 $\mu_3\int C_4$ 的 string-frame 规范中, RR 势的归一化不同, 不能把 $\mu_3$、此无量纲动能标准归一化荷以及张力直接当成三个相等量. 下一章用这些尺度建立 D3 的引力解和近视界字典.
